% problem_id: S001-theorem-flattening-local
% source_id: S001 (Stacks Project)
% source_file: flat.tex
% statement_locator: flat.tex lines 6789--6807
% proof_locator: flat.tex lines 6809--6902
% env: theorem
% id_note: label 在本来源内唯一
% status: complete (statement + author proof verbatim + reason + locators)
%
\begin{theorem}
\label{theorem-flattening-local}
In
Situation \ref{situation-flat-at-point}
assume $A$ is henselian, $B$ is essentially of finite type over $A$, and
$M$ is a finite $B$-module. Then there exists an ideal
$I \subset A$ such that $A/I$ corepresents the functor $F_{lf}$ on the category
$\mathcal{C}$. In other words given a local homomorphism of local rings
$\varphi : A \to A'$ with $B' = B \otimes_A A'$ and $M' = M \otimes_A A'$
the following are equivalent:
\begin{enumerate}
\item $\forall \mathfrak q \in V(\mathfrak m_{A'}B' + \mathfrak m_B B')
\subset \Spec(B') :
M'_{\mathfrak q}\text{ is flat over }A'$, and
\item $\varphi(I) = 0$.
\end{enumerate}
If $B$ is essentially of finite presentation over $A$ and $M$
of finite presentation over $B$, then $I$ is a finitely generated ideal.
\end{theorem}

\begin{proof}
Choose a finite type ring map $A \to C$ and a finite $C$-module $N$
and a prime $\mathfrak q$ of $C$ such that $B = C_{\mathfrak q}$
and $M = N_{\mathfrak q}$. In the following, when we say
``the theorem holds for $(N/C/A, \mathfrak q)$ we mean that
it holds for $(A \to B, M)$ where $B = C_{\mathfrak q}$ and
$M = N_{\mathfrak q}$. By
Lemma \ref{lemma-flat-at-point-go-up}
the functor $F_{lf}$ is unchanged if we replace $B$ by a local ring
flat over $B$. Hence, since $A$ is henselian, we may apply
Lemma \ref{lemma-existence-algebra}
and assume that there exists a complete d\'evissage of
$N/C/A$ at $\mathfrak q$.

\medskip\noindent
Let $(A_i, B_i, M_i, \alpha_i, \mathfrak q_i)_{i = 1, \ldots, n}$
be such a complete d\'evissage of $N/C/A$ at $\mathfrak q$. Let
$\mathfrak q'_i \subset A_i$ be the unique prime lying over
$\mathfrak q_i \subset B_i$ as in
Definition \ref{definition-complete-devissage-at-x-algebra}.
Since $C \to A_1$ is surjective and $N \cong M_1$ as $C$-modules,
we see by
Lemma \ref{lemma-flat-at-point-finite}
it suffices to prove the theorem holds for $(M_1/A_1/A, \mathfrak q'_1)$.
Since $B_1 \to A_1$ is finite and $\mathfrak q_1$ is the only prime
of $B_1$ over $\mathfrak q'_1$ we see that
$(A_1)_{\mathfrak q'_1} \to (B_1)_{\mathfrak q_1}$ is finite (see
Algebra, Lemma \ref{algebra-lemma-unique-prime-over-localize-below} or
More on Morphisms,
Lemma \ref{more-morphisms-lemma-finite-morphism-single-point-in-fibre}).
Hence by
Lemma \ref{lemma-flat-at-point-finite}
it suffices to prove the theorem holds for $(M_1/B_1/A, \mathfrak q_1)$.

\medskip\noindent
At this point we may assume, by induction on the length $n$ of the
d\'evissage, that the theorem holds for $(M_2/B_2/A, \mathfrak q_2)$.
(If $n = 1$, then $M_2 = 0$ which is flat over $A$.)
Reversing the last couple of steps of the previous paragraph, using
that $M_2 \cong \Coker(\alpha_2)$ as $B_1$-modules, we see
that the theorem holds for $(\Coker(\alpha_1)/B_1/A, \mathfrak q_1)$.

\medskip\noindent
Let $A'$ be an object of $\mathcal{C}$. At this point we use
Lemma \ref{lemma-induction-step}
to see that if $(M_1 \otimes_A A')_{\mathfrak q'}$ is flat
over $A'$ for a prime $\mathfrak q'$ of $B_1 \otimes_A A'$
lying over $\mathfrak m_{A'}$, then
$(\Coker(\alpha_1) \otimes_A A')_{\mathfrak q'}$ is flat over $A'$.
Hence we conclude that $F_{lf}$ is a subfunctor of the
functor $F'_{lf}$ associated to the module
$\Coker(\alpha_1)_{\mathfrak q_1}$ over $(B_1)_{\mathfrak q_1}$.
By the previous paragraph we know $F'_{lf}$ is corepresented by
$A/J$ for some ideal $J \subset A$. Hence we may replace $A$ by
$A/J$ and assume that $\Coker(\alpha_1)_{\mathfrak q_1}$ is
flat over $A$.

\medskip\noindent
Since $\Coker(\alpha_1)$ is a $B_1$-module for which
there exist a complete d\'evissage of $N_1/B_1/A$ at $\mathfrak q_1$
and since $\Coker(\alpha_1)_{\mathfrak q_1}$ is
flat over $A$ by
Lemma \ref{lemma-complete-devissage-flat-finite-type-module}
we see that $\Coker(\alpha_1)$ is free as an $A$-module, in particular
flat as an $A$-module. Hence
Lemma \ref{lemma-induction-step}
implies $F_{lf}(A')$ is nonempty if and only if $\alpha \otimes 1_{A'}$
is injective. Let $N_1 = \Im(\alpha_1) \subset M_1$ so that
we have exact sequences
$$
0 \to N_1 \to M_1 \to \Coker(\alpha_1) \to 0
\quad\text{and}\quad
B_1^{\oplus r_1} \to N_1 \to 0
$$
The flatness of $\Coker(\alpha_1)$ implies the first sequence
is universally exact (see
Algebra, Lemma \ref{algebra-lemma-flat-universally-injective}).
Hence $\alpha \otimes 1_{A'}$ is injective if and only if
$B_1^{\oplus r_1} \otimes_A A' \to N_1 \otimes_A A'$
is an isomorphism. Finally,
Theorem \ref{theorem-flattening-map}
applies to show this functor is corepresentable by $A/I$ for some ideal $I$
and we conclude $F_{lf}$ is corepresentable by $A/I$ also.

\medskip\noindent
To prove the final statement, suppose that $A \to B$ is essentially of finite
presentation and $M$ of finite presentation over $B$. Let $I \subset A$
be the ideal such that $F_{lf}$ is corepresented by $A/I$.
Write $I = \bigcup I_\lambda$ where $I_\lambda$ ranges over the finitely
generated ideals contained in $I$. Then, since $F_{lf}(A/I) = \{*\}$
we see that $F_{lf}(A/I_\lambda) = \{*\}$ for some $\lambda$, see
Lemma \ref{lemma-flat-at-point} part (2).
Clearly this implies that $I = I_\lambda$.
\end{proof}
